% GENERATED FILE. Edit canonical lessons, then run npm run generate:books.
% canonical-source-hash: fnv1a64:de586e3115bfec6e
% canonical-lessons: MR-C417-anamat, MR-C417-sthayik-hone, MR-C417-bhadekaru, MR-C417-tijori, MR-C417-udvahak

\chapter{Renting a Room}
\label{ch:mr-renting-a-room}
\hlchaptermodality{417}

\begin{chapteropening}
\textbf{By the end of this chapter:} \emph{I can talk about a deposit, settling in, a tenant, a strongbox and a lift.}

Complete the last lesson of chapter 417: I can talk about a deposit, settling in, a tenant, a strongbox and a lift.
\end{chapteropening}

\section[\texorpdfstring{\mr{अनामत} (anāmat)}{अनामत (anāmat)}]{\mr{अनामत} (anāmat) — a deposit; money held in trust}

\label{lesson:MR-C417-anamat}



\begin{quote}
Before the new one: say the Marathi for a porter, then the Marathi for to sprinkle.
\end{quote}

\subsection*{You'll want to know: \mr{अनामत}}
\textbf{\mr{अनामत}} — \emph{anāmat} — \textquotedblleft{}a deposit; money held in trust\textquotedblright{}.

\begin{quote}
\mr{खोली} \mr{सोडताना} \mr{अनामत} \mr{परत} \mr{मिळते}.
\end{quote}

\begin{quote}
\emph{kholī soḍtānā anāmat parat miḷte.}
\end{quote}

\begin{quote}
You get the deposit back when you leave the room.
\end{quote}

Feminine.

It comes back when you leave, if nothing is broken.

\begin{cousinweb}
\textbf{\mr{अनामत}} (\emph{anāmat}) is a loan through Persian, usually traced to Arabic \emph{amānat}, \textquotedblleft{}something held in trust, a deposit\textquotedblright{}.
\end{cousinweb}

\begin{grammarlens}[title={What you've built}]
One more word for getting around.
\end{grammarlens}

\subsection*{Your turn}
\begin{itemize}
\raggedright
  \item \emph{Say it:} \emph{anāmat}
  \item \emph{Say it:} \emph{anāmat}, once more
  \item \emph{Say it:} \emph{kholī soḍtānā anāmat parat miḷte.} — You get the deposit back when you leave the room.
  \item \emph{Recall:} say the Marathi for a porter, then the Marathi for to sprinkle, then say \emph{anāmat} again
\end{itemize}

\begin{tcolorbox}[breakable,colback=teal!4,colframe=teal!35!black,title={Before you move on}]
What does \mr{अनामत} mean? (\textquotedblleft{}A deposit; money held in trust\textquotedblright{}.) Say it once more.
\end{tcolorbox}
\section[\texorpdfstring{\mr{स्थायिक} \mr{होणे} (sthāyik hoṇe)}{स्थायिक होणे (sthāyik hoṇe)}]{\mr{स्थायिक} \mr{होणे} (sthāyik hoṇe) — to settle down; to make a home}

\label{lesson:MR-C417-sthayik-hone}



\begin{quote}
Before the new one: say the Marathi for to sprinkle, then the Marathi for a deposit.
\end{quote}

\subsection*{You'll want to know: \mr{स्थायिक} \mr{होणे}}
\textbf{\mr{स्थायिक} \mr{होणे}} — \emph{sthāyik hoṇe} — \textquotedblleft{}to settle down; to make a home\textquotedblright{}.

\begin{quote}
\mr{अरुणचं} \mr{कुटुंब} \mr{आता} \mr{पुण्यात} \mr{स्थायिक} \mr{झालं} \mr{आहे}.
\end{quote}

\begin{quote}
\emph{aruṇcaṁ kuṭumb ātā puṇyāt sthāyik jhālaṁ āhe.}
\end{quote}

\begin{quote}
Arun's family has now settled in Pune.
\end{quote}

Past (he …): \mr{तो} \mr{स्थायिक} \mr{झाला}, \emph{to sthāyik jhālā}.

The place takes -\mr{त}: \mr{पुण्यात} \mr{स्थायिक} \mr{होणे}, to settle in Pune for good.

\begin{cousinweb}
\textbf{\mr{स्थायिक} \mr{होणे}} (\emph{sthāyik hoṇe}) joins \textbf{\mr{स्थायिक}}, \textquotedblleft{}settled for good\textquotedblright{}, a tatsama from Sanskrit \emph{sthāyika}, \textquotedblleft{}lasting, enduring\textquotedblright{} (Monier-Williams), to the light verb \textbf{\mr{होणे}}, \textquotedblleft{}to become\textquotedblright{}.
\end{cousinweb}

\begin{grammarlens}[title={What you've built}]
One more word for getting around.
\end{grammarlens}

\subsection*{Your turn}
\begin{itemize}
\raggedright
  \item \emph{Say it:} \emph{sthāyik hoṇe}
  \item \emph{Say it:} \emph{sthāyik hoṇe}, once more
  \item \emph{Say it:} \emph{aruṇcaṁ kuṭumb ātā puṇyāt sthāyik jhālaṁ āhe.} — Arun's family has now settled in Pune.
  \item \emph{Recall:} say the Marathi for to sprinkle, then the Marathi for a deposit, then say \emph{sthāyik hoṇe} again
\end{itemize}

\begin{tcolorbox}[breakable,colback=teal!4,colframe=teal!35!black,title={Before you move on}]
What does \mr{स्थायिक} \mr{होणे} mean? (\textquotedblleft{}To settle down; to make a home\textquotedblright{}.) Say it once more.
\end{tcolorbox}
\section[\texorpdfstring{\mr{भाडेकरू} (bhāḍekarū)}{भाडेकरू (bhāḍekarū)}]{\mr{भाडेकरू} (bhāḍekarū) — a tenant}

\label{lesson:MR-C417-bhadekaru}



\begin{quote}
Before the new one: say the Marathi for a deposit, then the Marathi for to settle down.
\end{quote}

\subsection*{You'll want to know: \mr{भाडेकरू}}
\textbf{\mr{भाडेकरू}} — \emph{bhāḍekarū} — \textquotedblleft{}a tenant\textquotedblright{}.

\begin{quote}
\mr{शेजारच्या} \mr{घरात} \mr{नवे} \mr{भाडेकरू} \mr{आले}.
\end{quote}

\begin{quote}
\emph{śejārcyā gharāt nave bhāḍekarū āle.}
\end{quote}

\begin{quote}
New tenants have moved into the house next door.
\end{quote}

Masculine.

The owner is the \mr{मालक}; the plural is the same: \mr{दोन} \mr{भाडेकरू}.

\begin{cousinweb}
\textbf{\mr{भाडेकरू}} (\emph{bhāḍekarū}) is built from \textbf{\mr{भाडे}}, \textquotedblleft{}rent\textquotedblright{}, and the ending -\mr{करू}, \textquotedblleft{}one who does\textquotedblright{}: the one who pays the rent; the origin of \mr{भाडे} is not settled.
\end{cousinweb}

\begin{grammarlens}[title={What you've built}]
One more word for getting around.
\end{grammarlens}

\subsection*{Your turn}
\begin{itemize}
\raggedright
  \item \emph{Say it:} \emph{bhāḍekarū}
  \item \emph{Say it:} \emph{bhāḍekarū}, once more
  \item \emph{Say it:} \emph{śejārcyā gharāt nave bhāḍekarū āle.} — New tenants have moved into the house next door.
  \item \emph{Recall:} say the Marathi for a deposit, then the Marathi for to settle down, then say \emph{bhāḍekarū} again
\end{itemize}

\begin{tcolorbox}[breakable,colback=teal!4,colframe=teal!35!black,title={Before you move on}]
What does \mr{भाडेकरू} mean? (\textquotedblleft{}A tenant\textquotedblright{}.) Say it once more.
\end{tcolorbox}
\section[\texorpdfstring{\mr{तिजोरी} (tijorī)}{तिजोरी (tijorī)}]{\mr{तिजोरी} (tijorī) — a strongbox; a safe}

\label{lesson:MR-C417-tijori}



\begin{quote}
Before the new one: say the Marathi for to settle down, then the Marathi for a tenant.
\end{quote}

\subsection*{You'll want to know: \mr{तिजोरी}}
\textbf{\mr{तिजोरी}} — \emph{tijorī} — \textquotedblleft{}a strongbox; a safe\textquotedblright{}.

\begin{quote}
\mr{दुकानदार} \mr{पैसे} \mr{तिजोरीत} \mr{ठेवतो}.
\end{quote}

\begin{quote}
\emph{dukāndār paise tijorīt ṭhevto.}
\end{quote}

\begin{quote}
The shopkeeper keeps the money in the safe.
\end{quote}

Feminine.

Also the treasury of a town or a state: \mr{सरकारी} \mr{तिजोरी}.

\begin{cousinweb}
\textbf{\mr{तिजोरी}} (\emph{tijorī}) is a Marathi word for a steel box with a lock, and for a treasury; where it comes from is not settled.
\end{cousinweb}

\begin{grammarlens}[title={What you've built}]
One more word for getting around.
\end{grammarlens}

\subsection*{Your turn}
\begin{itemize}
\raggedright
  \item \emph{Say it:} \emph{tijorī}
  \item \emph{Say it:} \emph{tijorī}, once more
  \item \emph{Say it:} \emph{dukāndār paise tijorīt ṭhevto.} — The shopkeeper keeps the money in the safe.
  \item \emph{Recall:} say the Marathi for to settle down, then the Marathi for a tenant, then say \emph{tijorī} again
\end{itemize}

\begin{tcolorbox}[breakable,colback=teal!4,colframe=teal!35!black,title={Before you move on}]
What does \mr{तिजोरी} mean? (\textquotedblleft{}A strongbox; a safe\textquotedblright{}.) Say it once more.
\end{tcolorbox}
\section[\texorpdfstring{\mr{उद्वाहक} (udvāhak)}{उद्वाहक (udvāhak)}]{\mr{उद्वाहक} (udvāhak) — an elevator; a lift}

\label{lesson:MR-C417-udvahak}



\begin{quote}
Before the new one: say the Marathi for a tenant, then the Marathi for a strongbox.
\end{quote}

\subsection*{You'll want to know: \mr{उद्वाहक}}
\textbf{\mr{उद्वाहक}} — \emph{udvāhak} — \textquotedblleft{}an elevator; a lift\textquotedblright{}.

\begin{quote}
\mr{इमारतीतला} \mr{उद्वाहक} \mr{बंद} \mr{आहे}.
\end{quote}

\begin{quote}
\emph{imāratītlā udvāhak band āhe.}
\end{quote}

\begin{quote}
The lift in the building is not working.
\end{quote}

Masculine.

A formal word for signs; in speech many people say \mr{लिफ्ट}.

\begin{cousinweb}
\textbf{\mr{उद्वाहक}} (\emph{udvāhak}) is a modern word built from Sanskrit \emph{ud-}, \textquotedblleft{}up\textquotedblright{}, and \emph{vāhaka}, \textquotedblleft{}a carrier\textquotedblright{} (Monier-Williams): a thing that carries up.
\end{cousinweb}

\begin{grammarlens}[title={What you've built}]
One more word for getting around.
\end{grammarlens}

\subsection*{Your turn}
\begin{itemize}
\raggedright
  \item \emph{Say it:} \emph{udvāhak}
  \item \emph{Say it:} \emph{udvāhak}, once more
  \item \emph{Say it:} \emph{imāratītlā udvāhak band āhe.} — The lift in the building is not working.
  \item \emph{Recall:} say the Marathi for a tenant, then the Marathi for a strongbox, then say \emph{udvāhak} again
\end{itemize}

\begin{tcolorbox}[breakable,colback=teal!4,colframe=teal!35!black,title={Before you move on}]
What does \mr{उद्वाहक} mean? (\textquotedblleft{}An elevator; a lift\textquotedblright{}.) Say it once more.
\end{tcolorbox}
